\documentclass[11pt]{article}
\usepackage{amsmath,amssymb,amsthm,amscd}
\usepackage[margin=1in]{geometry}
\usepackage{graphicx}
\usepackage[colorlinks=true,linkcolor=blue,urlcolor=blue]{hyperref}

\newtheorem{definition}{Definition}
\newtheorem{proposition}{Proposition}
\newtheorem{remark}{Remark}

\newcommand{\M}{\mathcal{M}}
\newcommand{\B}{\mathcal{B}}
\newcommand{\Hh}{\mathcal{H}}
\newcommand{\T}{T}
\newcommand{\R}{\mathbb{R}}
\newcommand{\Sym}{\mathrm{Sym}}
\newcommand{\tr}{\operatorname{tr}}
\newcommand{\rank}{\operatorname{rank}}
\newcommand{\dd}{\mathrm{d}}

\title{\bf Pullback Metrics for Neural Computation:\\ Mathematical Formulation}
\author{PullbackMetric project}
\date{\today}

\begin{document}
\maketitle

\noindent
This note states, in precise geometric terms, every object the pipeline manipulates: the
manifolds, the maps between them, the metric on the output manifold, the pullback metric it
induces on neural--state space, and the quantities we read off. Each definition is followed
by what we actually compute. Output figures appear only at the end (\S\ref{sec:results}).

\section{Objects: manifolds and maps}\label{sec:objects}

\begin{definition}[Neural--state manifold]
Let $\M$ be the \emph{neural--state manifold}: a smooth manifold of dimension $k$ whose
points are population activity states. In practice $\M$ is an open subset of $\R^{k}$ with a
single global chart $z=(z^{1},\dots,z^{k})$ (the top-$k$ principal-component scores of the
binned firing rate; see \S\ref{sec:compute}). Its tangent space at $z$ is
$\T_{z}\M\cong\R^{k}$; a tangent vector $v\in \T_z\M$ is an infinitesimal change of neural
state.
\end{definition}

\begin{definition}[Behavior manifold]
Let $\B$ be the \emph{behavior manifold}: a smooth manifold of dimension $m$ whose points are
behavioral outputs. For single-trial lick timing, $\B=(0,\infty)\subset\R$ (a positive
scalar), so $m=1$; in general $\B\subseteq\R^{m}$ (e.g.\ $[\text{lick time},\ P(\text{reward})]$,
$m=2$, or a $c$-way choice, $m=c$). Coordinates on $\B$ are $y=(y^{1},\dots,y^{m})$.
\end{definition}

\begin{definition}[Readout map]
The \emph{readout} is a smooth map
\[
  f:\M\longrightarrow\B ,\qquad z\mapsto y=f(z),
\]
the trained model that produces behavior from neural state. When the model is layered we
factor it through a hidden manifold $\Hh\cong\R^{p}$,
\[
  \begin{CD}
    \M @>{f_{1}}>> \Hh @>{f_{2}}>> \B ,
  \end{CD}
  \qquad f=f_{2}\circ f_{1}.
\]
\end{definition}

\begin{definition}[Differential / Jacobian]
The differential of $f$ at $z$ is the linear map between tangent spaces
\[
  \dd f_{z}:\ \T_{z}\M\longrightarrow \T_{f(z)}\B ,
\]
represented in the coordinates $z,y$ by the Jacobian matrix
\[
  J(z)\in\R^{m\times k},\qquad J^{a}{}_{i}(z)=\frac{\partial f^{a}}{\partial z^{i}}(z),
  \qquad (\dd f_{z}\,v)^{a}=J^{a}{}_{i}\,v^{i}.
\]
For a layered readout the chain rule gives $\dd f_{z}=(\dd f_{2})_{f_1(z)}\circ(\dd f_{1})_{z}$,
i.e.\ $J(z)=J_{2}\,J_{1}$.
\end{definition}

\section{A Riemannian metric on the output manifold}\label{sec:outmetric}

\begin{definition}[Riemannian metric]
A \emph{Riemannian metric} $h$ on $\B$ is a smooth field of symmetric, positive-definite
bilinear forms
\[
  h_{y}:\ \T_{y}\B\times \T_{y}\B\to\R ,\qquad
  h = h_{ab}(y)\,\dd y^{a}\otimes \dd y^{b},
\]
with $h_{ab}=h_{ba}$ and $(h_{ab})\succ0$. It measures lengths on $\B$: a curve
$\beta:[0,1]\to\B$ has length $L_{h}[\beta]=\int_{0}^{1}\!\sqrt{h_{\beta(t)}(\dot\beta,\dot\beta)}\,\dd t$,
and $d_{\B}(y_{0},y_{1})=\inf_{\beta}L_{h}[\beta]$.
\end{definition}

The choice of $h$ encodes \emph{what ``distance between behaviors'' means}. We take $h$ to be
the Fisher--Rao metric of an observation model for the behavior --- the unique (up to scale)
metric invariant under reparametrization (Chentsov). Concretely:

\begin{itemize}
\item \textbf{Euclidean:} $h_{ab}=\delta_{ab}$ (raw output units).
\item \textbf{Gaussian} $Y\sim\mathcal N(y,\sigma^{2})$ with fixed $\sigma$: the Fisher metric
      in the mean is $h_{ab}=\sigma^{-2}\delta_{ab}$; for $m=1$, $h=\sigma^{-2}\,\dd y\otimes\dd y$.
      Distance counts standard deviations of behavioral change.
\item \textbf{Logarithmic (Weber)} model $\log Y\sim\mathcal N(\mu,\sigma^{2})$ parametrized by
      the median $y=e^{\mu}$: by change of variables $\mu=\log y$,
      \[
        h_{yy}(y)=\Big(\tfrac{\dd\mu}{\dd y}\Big)^{2}\frac1{\sigma^{2}}
                 =\frac{1}{\sigma^{2}y^{2}},\qquad
        h=\frac{1}{\sigma^{2}}\,(\dd\log y)^{2}.
      \]
      A fixed \emph{fractional} change in $y$ has constant length --- the scale appropriate to
      interval timing (scalar variability).
\item \textbf{Bernoulli} $Y\sim\mathrm{Bern}(y)$: $h_{yy}=\big(y(1-y)\big)^{-1}$;
      \textbf{categorical} (softmax logits $y$): $h=\mathrm{diag}(p)-pp^{\top}$, $p=\mathrm{softmax}(y)$.
\end{itemize}

\section{The pullback metric}\label{sec:pullback}

\begin{definition}[Pullback metric]
Given the smooth readout $f:\M\to\B$ and a metric $h$ on $\B$, the \emph{pullback} $g=f^{*}h$
is the symmetric $(0,2)$ tensor field on $\M$ defined pointwise by
\[
  g_{z}(u,v)\ :=\ h_{f(z)}\big(\dd f_{z}\,u,\ \dd f_{z}\,v\big),
  \qquad u,v\in \T_{z}\M .
\]
In coordinates,
\[
  \boxed{\,g_{ij}(z)=h_{ab}\big(f(z)\big)\,\frac{\partial f^{a}}{\partial z^{i}}\,
                     \frac{\partial f^{b}}{\partial z^{j}}
        \quad\Longleftrightarrow\quad
        g(z)=J(z)^{\top}\,H\big(f(z)\big)\,J(z)\,,}
\]
where $H=(h_{ab})$ is the output-metric matrix.
\end{definition}

\begin{proposition}[Basic properties]\label{prop:props}
Let $g=f^{*}h$ with $H\succ0$. Then
\begin{enumerate}
\item[(i)] $g$ is symmetric and positive \emph{semi}definite: $v^{\top}g\,v=
      h(\dd f\,v,\dd f\,v)\ge 0$, with equality iff $\dd f_{z}\,v=0$.
\item[(ii)] $\ker g_{z}=\ker \dd f_{z}$. Hence $g_{z}$ is a genuine (positive-definite)
      Riemannian metric iff $\dd f_{z}$ is injective, i.e.\ $\rank J(z)=k$ (requires $k\le m$);
      otherwise $g$ is a \emph{degenerate} pseudometric.
\item[(iii)] The null space $\ker\dd f_{z}$ consists of neural-state directions that leave the
      behavior unchanged to first order --- the infinitesimal \emph{equivalence classes} the
      computation collapses.
\item[(iv)] Pullback is functorial: $(f_{2}\circ f_{1})^{*}h=f_{1}^{*}(f_{2}^{*}h)$, i.e.\
      $g=J_{1}^{\top}J_{2}^{\top}H\,J_{2}J_{1}$.
\end{enumerate}
\end{proposition}

\begin{proposition}[Isometric meaning of length]\label{prop:length}
For any curve $\gamma$ in $\M$, $L_{g}[\gamma]=L_{h}[f\circ\gamma]$: the pullback length of a
state path equals the behavioral length of its image. Consequently
$d_{\M}(z_{0},z_{1})\ge d_{\B}\big(f(z_{0}),f(z_{1})\big)$, with equality when the image of the
$g$-geodesic is an $h$-geodesic. Thus pullback distances measure \emph{behavioral
distinguishability of neural states}, not raw neural distance.
\end{proposition}

\paragraph{What $g$ looks like.}
\emph{Scalar output} ($m=1$), $h=\rho(y)\,\dd y\otimes\dd y$ with $\rho>0$. Writing the
gradient one-form $w_{i}=\partial_{i}f$ (so $\dd f=w_{i}\,\dd z^{i}$),
\[
  g_{ij}(z)=\rho\big(f(z)\big)\,w_{i}w_{j},\qquad
  g=\rho(f)\,\dd f\otimes \dd f .
\]
This is \textbf{rank one} at every $z$: one non-null direction (the gradient $w$) and a
$(k-1)$-dimensional kernel $w^{\perp}$. The magnitude of $g$ is set by $\rho(f(z))$:
\[
  \text{Gaussian: } \rho=\tfrac1{\sigma^{2}}\ \Rightarrow\ g=\tfrac1{\sigma^{2}}\,\dd f\otimes\dd f,
  \qquad
  \text{log: } \rho=\tfrac1{\sigma^{2}f^{2}}\ \Rightarrow\ g=\tfrac1{\sigma^{2}f^{2}}\,\dd f\otimes\dd f .
\]
\emph{Vector output} ($m\ge2$): $\rank g_{z}\le\min(m,k)$, generically $m$, so $g$ has an
$m$-dimensional non-null subspace and the local unit balls are ellipsoids.

\begin{remark}[Enlarging the network does not raise the rank]
For a scalar readout, $\rank g_{z}=1$ regardless of the size or depth of $f$. A larger
network changes how the single non-null direction $w(z)$ and the scale $\rho(f(z))$
\emph{vary} across $\M$ (curvature / state-dependence), not the pointwise rank.
\end{remark}

\section{Derived geometric quantities}\label{sec:derived}

\paragraph{Spectrum requires a reference metric.}
Diagonalizing a $(0,2)$ tensor is only meaningful relative to a reference inner product
$\gamma$ on $\T_{z}\M$ (used to identify vectors with covectors). We take $\gamma$ to be the
Euclidean metric in the PC chart, $\gamma=I_{k}$. The \emph{principal directions} solve the
generalized eigenproblem
\[
  g\,v=\lambda\,\gamma\,v \quad(\gamma=I:\ g\,v=\lambda v),\qquad
  \lambda_{\max}=\max_{v\ne0}\frac{v^{\top}g\,v}{v^{\top}\gamma\,v}
              =\max_{v\ne0}\frac{h(\dd f\,v,\dd f\,v)}{\|v\|^{2}} .
\]
So $\lambda_{\max}$ is the largest squared behavioral change per unit (reference) step, its
eigenvector is the \emph{most behavior-sensitive} state direction, and
$\sqrt{\lambda_{\max}}$ is the corresponding \emph{sensitivity} (dimensionless: units of
$h$ per unit reference length). For the scalar case, $\lambda_{\max}=\rho(f)\|w\|^{2}$ and
$\sqrt{\lambda_{\max}}=\sqrt{\rho(f)}\,\|w\|$, with eigenvector $w/\|w\|$. This spectrum is
basis-dependent; a different chart or reference $\gamma$ (e.g.\ the neural covariance) changes
it.

\paragraph{Volume, rank.}
Where $g\succ0$, the Riemannian volume element is $\dd V_{g}=\sqrt{\det g}\ \dd z^{1}\!\wedge\!\cdots\!\wedge\dd z^{k}$;
for degenerate $g$ we use the pseudo-volume $\prod_{\lambda_{i}>0}\sqrt{\lambda_{i}}$. The
effective rank $\exp\!\big(-\sum_i p_i\log p_i\big)$ with $p_i=\lambda_i/\sum_j\lambda_j$, and
the participation ratio $(\sum_i\lambda_i)^2/\sum_i\lambda_i^2$, count how many directions
matter locally.

\paragraph{Connection and curvature.}
Where $g$ is nondegenerate it has a Levi-Civita connection with Christoffel symbols
\[
  \Gamma^{k}{}_{ij}=\tfrac12 g^{kl}\big(\partial_{i}g_{jl}+\partial_{j}g_{il}-\partial_{l}g_{ij}\big),
\]
geodesics $\ddot\gamma^{k}+\Gamma^{k}{}_{ij}\dot\gamma^{i}\dot\gamma^{j}=0$, Riemann tensor
$R^{\rho}{}_{\sigma\mu\nu}=\partial_{\mu}\Gamma^{\rho}{}_{\nu\sigma}-\partial_{\nu}\Gamma^{\rho}{}_{\mu\sigma}
+\Gamma^{\rho}{}_{\mu\lambda}\Gamma^{\lambda}{}_{\nu\sigma}-\Gamma^{\rho}{}_{\nu\lambda}\Gamma^{\lambda}{}_{\mu\sigma}$,
Ricci $R_{\sigma\nu}=R^{\rho}{}_{\sigma\rho\nu}$ and scalar $R=g^{\sigma\nu}R_{\sigma\nu}$.

\section{What we compute, object by object}\label{sec:compute}

\begin{enumerate}
\item \textbf{The chart $\M$ (neural state).} From spikes we bin at $20$\,ms, average the rate
      over an analysis window (e.g.\ $[0,200]$\,ms post-cue), standardize per unit
      $s=(r-\mu)\oslash\varsigma$, and project onto the top-$k$ PC loadings $U\in\R^{N\times k}$:
      $z=U^{\top}s$. This composite is the coordinate map defining $\M$ ($k=6$ here).
      \emph{(modules: \texttt{data.load\_session}, \texttt{features.window\_features},
      \texttt{reduce.PCAState}.)}
\item \textbf{The readout $f$.} Fit to predict behavior from $z$: ridge (linear,
      $f(z)=w^{\top}z+b$) or a small MLP ($f(z)=A_{2}\tanh(A_{1}z+b_{1})+b_{2}$).
      $\sigma$ is the cross-validated residual std (and $\sigma_{\log}$ the residual std of
      $\log y$). \emph{(\texttt{maps.ReadoutMap} $+$ scikit-learn.)}
\item \textbf{The Jacobian $J$.} For the ridge readout $J=w^{\top}$ (constant). For the MLP
      $f(z)=A_{2}\,\phi(A_{1}z+b_{1})+b_{2}$ (activation $\phi$) the exact Jacobian is
      \[
        J(z)=A_{2}\,\mathrm{diag}\!\big(\phi'(A_{1}z+b_{1})\big)\,A_{1}
        \quad(\phi=\tanh:\ \phi'=1-\tanh^{2}),
      \]
      and for $L$ layers $J=W_{L}^{\top}D_{L-1}W_{L-1}^{\top}\cdots D_{1}W_{1}^{\top}$ with
      $D_{l}=\mathrm{diag}(\phi'(a_{l}))$ evaluated at the pre-activations $a_l$; a leading
      linear map (standardizer, PCA) contributes a constant factor $P$ by the chain rule,
      $J_{\text{tot}}=J_{\mathrm{mlp}}(Pz+o)\,P$. The demo figures were produced with the
      central finite-difference Jacobian
      $J^{a}{}_{i}\approx\big(f^{a}(z+\delta e_{i})-f^{a}(z-\delta e_{i})\big)/(2\delta)$
      (\texttt{maps.numeric\_jacobian}); the closed form above is implemented in
      \texttt{maps.MLPReadout} and agrees with the finite-difference Jacobian to
      $\sim\!10^{-11}$ on the fitted network. Because $\dim\B=1$, $J$ is a single row
      $\nabla f(z)^{\top}$ either way, and $g$ is rank one (Rmk.~1).
\item \textbf{The output metric $H$.} The closed-form matrices of \S\ref{sec:outmetric}.
      \emph{(\texttt{output\_metrics}.)}
\item \textbf{The pullback $g=J^{\top}HJ$}, formed densely, or applied matrix-free as
      $g\,v=J^{\top}\!\big(H(Jv)\big)$. \emph{(\texttt{pullback.PullbackMetric}.)}
\item \textbf{Spectrum} by symmetric eigendecomposition of $g$ (reference $\gamma=I$);
      volume, effective rank, participation ratio as above.
\item \textbf{Connection/curvature} by finite differencing $g$ (\texttt{geometry.RiemannianField}).
      Validated on analytic cases: $g=J^{\top}J$ reproduces $W^{\top}W$ exactly; an isometric
      $f$ gives $g=I$; the round $2$-sphere gives scalar curvature $2/r^{2}$ to $<2\%$; a flat
      metric gives straight-line geodesics.
\end{enumerate}

\section{Results}\label{sec:results}

\begin{figure}[h]
\centering
\includegraphics[width=\textwidth]{../figures/demo_log_metric.png}
\caption{Same neural state ($k=6$ PCs, SM318) and the same linear readout $f$; only the
output metric $h$ differs. \textbf{Left:} sensitivity $\sqrt{\lambda_{\max}}$ vs.\ predicted
lick time $f(z)$. Gaussian $h=\sigma^{-2}$ gives $\sqrt{\lambda_{\max}}=\|w\|/\sigma$,
constant (flat red line), because a linear readout has one slope everywhere. Logarithmic
$h=(\sigma^{2}y^{2})^{-1}$ gives $\sqrt{\lambda_{\max}}=\|w\|/(\sigma_{\log}f)$, the declining
$1/f$ curve. \textbf{Right:} the log-metric sensitivity over the PC$_1$--PC$_2$ plane (own
color scale); bright where predicted licks are early, contours are level sets of $f$, arrows
are the sensitive direction $w/\|w\|$.}
\end{figure}

\begin{figure}[h]
\centering
\includegraphics[width=\textwidth]{../figures/demo_metric_visualization.png}
\caption{Visualizing $g$ on the top two PCs. \textbf{Left (scalar readout, rank-1 metric):}
background $\sqrt{\lambda_{\max}}$, white contours of $f$ (iso-metric level sets), arrows the
single behavior-sensitive direction. \textbf{Right (vector readout $[\,\text{lick},P(\text{reward})\,]$,
rank-2):} the sensitivity-ellipse field (image of the unit circle under $g^{1/2}$; long axis =
most behavior-sensitive direction), whose orientation varies across state space.}
\end{figure}

\noindent
Interpretation and the broader plan are in \texttt{docs/METHOD.md} and
\texttt{docs/pullback\_metric\_survey\_and\_roadmap.md}.

\end{document}
